
\documentclass{article}
\usepackage{colortbl}
\usepackage{makecell}
\usepackage{multirow}
\usepackage{supertabular}

\begin{document}

\newcounter{utterance}

\centering \large Interaction Transcript for game `taboo', experiment `low\_en', episode 16 with reflexion{-}Qwen3.5{-}9B.
\vspace{24pt}

{ \footnotesize  \setcounter{utterance}{1}
\setlength{\tabcolsep}{0pt}
\begin{supertabular}{c@{$\;$}|p{.15\linewidth}@{}p{.15\linewidth}p{.15\linewidth}p{.15\linewidth}p{.15\linewidth}p{.15\linewidth}}
    \# & $\;$A & \multicolumn{4}{c}{Game Master} & $\;\:$B\\
    \hline

    \theutterance \stepcounter{utterance}  
    & & \multicolumn{4}{p{0.6\linewidth}}{
        \cellcolor[rgb]{0.9,0.9,0.9}{
            \makecell[{{p{\linewidth}}}]{
                \texttt{\tiny{[P1$\langle$GM]}}
                \texttt{You have attempted this same task before and did not succeed. These are the plans you wrote after those attempts.} \\
\texttt{Trial \#1: New plan: I failed because my clue explicitly mentioned "induce sleep," which directly used the related word "sleep" from the forbidden list. Next time, I will describe the drug's mechanism (e.g., "a central nervous system depressant used for anxiety") or its chemical class ("a benzodiazepine") without naming any of the provided related words like sleep, narcotic, or diazepam. I will ensure the description focuses solely on the brand name "Valium" and its general medical use without referencing the specific synonyms or categories listed in the rules.} \\
\\ 
\\ 
\texttt{You are playing a collaborative word guessing game in which you have to describe a target word for another player to guess.} \\
\\ 
\texttt{Rules:} \\
\texttt{(a) You have to reply in the form: CLUE: <some text>. Guesses from the other player will start with GUESS.} \\
\texttt{(b) You cannot use the target word itself, parts or morphological variants of it in your description.} \\
\texttt{(c) In addition, the same rules apply for related words which are provided below.} \\
\\ 
\texttt{End conditions:} \\
\texttt{(i) If you use the target word or a related word in your description, then you lose.} \\
\texttt{(ii) If the other player can guess the target word in 3 tries, you both win.} \\
\\ 
\texttt{Let us start.} \\
\\ 
\texttt{This is the target word that you need to describe and that the other player needs to guess:} \\
\\ 
\texttt{valium} \\
\\ 
\texttt{Related words are:} \\
\\ 
\texttt{{-} sleep} \\
\texttt{{-} narcotic} \\
\texttt{{-} diazepam} \\
\\ 
\texttt{Important: You are under time pressure, give short descriptions that are to the point!} \\
            }
        }
    }
    & \\ \\

    \theutterance \stepcounter{utterance}  
    & \multicolumn{4}{p{0.6\linewidth}}{
        \cellcolor[rgb]{0.9,0.9,0.9}{
            \makecell[{{p{\linewidth}}}]{
                \texttt{\tiny{[P1$\rangle$GM]}}
                \texttt{CLUE: A benzodiazepine brand name often prescribed for anxiety and muscle spasms.} \\
            }
        }
    }
    & & \\ \\

    \theutterance \stepcounter{utterance}  
    & & & \multicolumn{2}{p{0.3\linewidth}}{
        \cellcolor[rgb]{0.9,0.9,0.9}{
            \makecell[{{p{\linewidth}}}]{
                \texttt{\tiny{[GM$|$GM]}}
                \texttt{continue} \\
            }
        }
    }
    & & \\ \\

    \theutterance \stepcounter{utterance}  
    & & & \multicolumn{2}{p{0.3\linewidth}}{
        \cellcolor[rgb]{0.9,0.9,0.9}{
            \makecell[{{p{\linewidth}}}]{
                \texttt{\tiny{[GM$|$GM]}}
                \texttt{a benzodiazepine brand name often prescribed for anxiety and muscle spasms} \\
            }
        }
    }
    & & \\ \\

    \theutterance \stepcounter{utterance}  
    & & \multicolumn{4}{p{0.6\linewidth}}{
        \cellcolor[rgb]{0.9,0.9,0.9}{
            \makecell[{{p{\linewidth}}}]{
                \texttt{\tiny{[GM$\rangle$P2]}}
                \texttt{You are playing a collaborative word guessing game in which you have to guess a target word that another player describes to you.} \\
\\ 
\texttt{You can make one guess at each trial. You win when you guess the target word. You lose when you cannot guess it in 3 tries.} \\
\\ 
\texttt{After each trial you will get a new hint from the other player which starts with CLUE.} \\
\\ 
\texttt{Make your guesses by just saying the word using the following form: GUESS: <a word>} \\
\\ 
\texttt{Let us start.} \\
\\ 
\\ 
\texttt{CLUE: A benzodiazepine brand name often prescribed for anxiety and muscle spasms.} \\
            }
        }
    }
    & \\ \\

    \theutterance \stepcounter{utterance}  
    & & & \multicolumn{4}{p{0.6\linewidth}}{
        \cellcolor[rgb]{0.9,0.9,0.9}{
            \makecell[{{p{\linewidth}}}]{
                \texttt{\tiny{[GM$\langle$P2]}}
                \texttt{GUESS: Valium} \\
            }
        }
    }
     \\ \\

    \theutterance \stepcounter{utterance}  
    & & & \multicolumn{2}{p{0.3\linewidth}}{
        \cellcolor[rgb]{0.9,0.9,0.9}{
            \makecell[{{p{\linewidth}}}]{
                \texttt{\tiny{[GM$|$GM]}}
                \texttt{continue} \\
            }
        }
    }
    & & \\ \\

    \theutterance \stepcounter{utterance}  
    & & & \multicolumn{2}{p{0.3\linewidth}}{
        \cellcolor[rgb]{0.9,0.9,0.9}{
            \makecell[{{p{\linewidth}}}]{
                \texttt{\tiny{[GM$|$GM]}}
                \texttt{valium} \\
            }
        }
    }
    & & \\ \\

    \theutterance \stepcounter{utterance}  
    & & & \multicolumn{2}{p{0.3\linewidth}}{
        \cellcolor[rgb]{0.9,0.9,0.9}{
            \makecell[{{p{\linewidth}}}]{
                \texttt{\tiny{[GM$|$GM]}}
                \texttt{end game} \\
            }
        }
    }
    & & \\ \\

\end{supertabular}
}

\end{document}
